\chapter{Convenció de mida}
\label{cap:mida}

Al llarg del llibre hem parlat de categories \emph{petites}, \emph{localment petites} i de \emph{classes pròpies} sense fixar del tot el marc conjuntista. Aquest apèndix el fa explícit, perquè el volum de lògica categòrica ---on apareixen categories de models, doctrines i fibracions--- l'exigirà.

\medskip
\noindent\textbf{Marc adoptat: classes i conjunts.} Treballem en una teoria de conjunts amb \textbf{classes pròpies}, a l'estil de von Neumann--Bernays--Gödel (NBG). Hi ha dos nivells de col·leccions: els \emph{conjunts}, que poden ser elements d'altres col·leccions, i les \emph{classes pròpies}, que no. Amb això:
\begin{itemize}
  \item una categoria és \textbf{petita} si la seva col·lecció d'objectes i la de morfismes són totes dues \emph{conjunts};
  \item és \textbf{localment petita} si tots els homconjunts $\Hom(A,B)$ són conjunts, tot i que la col·lecció d'objectes pot ser una classe pròpia;
  \item és \textbf{ben potenciada} si, per a cada objecte $A$, els subobjectes de $A$ formen un conjunt.
\end{itemize}
L'última condició demana una precisió. Un \emph{subobjecte} de $A$ no és un monomorfisme concret cap a $A$, sinó una classe d'equivalència de monomorfismes. Dos monomorfismes $m\colon S\rightarrowtail A$ i $n\colon T\rightarrowtail A$ són equivalents quan cadascun factoritza a través de l'altre, és a dir, quan difereixen en un isomorfisme entre els dominis. La distinció és essencial, i en el marc de classes cal formular-la amb cura. A $\cat{Set}$, els monomorfismes cap a un conjunt donat formen una classe pròpia, perquè cada subconjunt té moltes còpies isomorfes; i cada classe d'equivalència és, ella mateixa, una classe pròpia. Com que les classes pròpies no poden ser elements de cap col·lecció, \enquote{el conjunt de les classes d'equivalència} no té sentit literal. La definició precisa és, doncs, la següent: una categoria és \textbf{ben potenciada} si, per a cada objecte $A$, hi ha un \emph{conjunt} $R_A$ de monomorfismes cap a $A$ tal que tot monomorfisme cap a $A$ és equivalent a algun element de $R_A$. Aleshores es defineix $\Sub(A)$ com el conjunt quocient $R_A/{\sim}$, i $[m]$ designa la classe dels representants equivalents a $m$. Un altre conjunt de representants dona un quocient en bijecció canònica amb aquest, i la bijecció respecta l'ordre; per això el resultat no depèn de la tria de $R_A$. A $\cat{Set}$ es pot prendre com a $R_A$ el conjunt de les inclusions dels subconjunts de $A$, i s'obté $\Sub(A)\cong\mathcal P(A)$. La resta del llibre parla de \enquote{classes de monomorfismes} en aquest sentit.

Les categories $\cat{Set}$, $\cat{Grp}$, $\cat{Top}$, $\cat{Graph}$ i $\cat{Cat}$ són localment petites però no petites; cada preordre concret sobre un conjunt i cada categoria finita concreta són petites. Això no s'ha de confondre amb la petitesa de la categoria de \emph{tots} els preordres o de \emph{totes} les categories finites, que ---en tenir una classe pròpia de conjunts subjacents--- són grans llevat que es fixi una convenció de mida o s'esculli un conjunt de representants.

Quan alguna construcció demani triar simultàniament un objecte (o un morfisme) per a cada objecte d'una categoria \emph{gran} ---per exemple, en construir un invers d'una equivalència---, la família de tries està indexada per una classe pròpia, i l'axioma de l'elecció ordinari, que val per a famílies indexades per un conjunt, no hi arriba. En aquests casos assumim la forma \textbf{global} de l'elecció (disponible, per exemple, a NBG amb l'axioma d'elecció global). Ho indicarem allà on calgui.

\medskip
\noindent\textbf{L'alternativa dels universos.} Una segona opció, molt usada en geometria algebraica i teoria de topos, fixa un o més \textbf{universos de Grothendieck} $\mathcal{U}$ ---conjunts prou grans i tancats per les operacions habituals--- i parla de categories $\mathcal{U}$-petites i $\mathcal{U}$-grans de manera \emph{relativa} a $\mathcal{U}$. Aquesta convenció té l'avantatge que \enquote{la categoria de tots els $\mathcal{U}$-conjunts} és, ella mateixa, un objecte d'un univers més gran, cosa que simplifica certes construccions. En aquest volum \emph{no} l'adoptem, per no carregar una introducció amb la teoria d'universos; però el lector que continuï cap a la teoria de feixos la retrobarà. Les dues opcions són maneres alternatives de \emph{gestionar la mida} més que no pas fonaments idèntics: a grans trets, el paper de les \enquote{classes pròpies} de NBG el fa aquí la distinció entre $\mathcal{U}$-petit i $\mathcal{U}$-gran relativa a un univers $\mathcal{U}$ fixat. Fer precisa aquesta correspondència exigeix especificar l'univers i les hipòtesis pertinents; no cal, però, que això sigui un obstacle per a una primera lectura.

\medskip
\noindent Fixar aquesta convenció des d'ara evita l'ambigüitat que, altrament, es faria notar quan apareguin categories de models i functors classificadors: en tot moment, cada construcció d'aquest llibre declara si viu en un conjunt o en una classe pròpia.

\medskip
\noindent\textbf{Per saber-ne més.} Les qüestions de mida es tracten, al nivell d'aquest llibre, a \cite[§1.8]{awodey}, \cite[§1.1]{riehl} ---on l'observació~1.1.5 presenta la convenció d'universos---, \cite[§3.2]{leinster} i \cite[cap.~1]{perrone}. El capítol~3 de \cite{leinster} presenta, a més, la perspectiva en què el mateix llenguatge estructural fa de fonament, sense una teoria de conjunts prèvia; en són les fonts \cite{lawvere-etcs} i, com a introducció informal, \cite{leinster-rethinking}. Per a una comparació detallada de les diferents maneres de gestionar la mida ---classes, universos i altres alternatives---, adequada un cop se'n té una certa experiència, vegeu \cite{shulman-mida}.
